\documentclass[12pt]{article}
\usepackage{graphicx}
\usepackage{mathtools,amsfonts,amsthm}
\usepackage{lipsum}
\usepackage{enumitem}% http://ctan.org/pkg/enumitem

\begin{document}

\section*{LITERATURE REVIEW}
In this chapter, a review of the existing research and scholarship related to mental health and its effects on a person’s productivity as well as the interaction between an individual’s cognitive, affective and somatic symptoms of depression is made. The following are some of the previous works.\\

\noindent A research study conducted by Martin et al. (2009), the authors investigated whether different types of health promotion interventions in the workplace reduced depression and anxiety symptoms. A systematic review and meta-analysis of the literature was undertaken on workplace health promotion published during the period 1997-2007.Studies were considered eligible for inclusion if they evaluated the impact of an intervention using a valid indicator or specific measure of depression or anxiety symptoms. The standardized mean difference was calculated for each of the following three types of outcome measures; depression, anxiety and composite mental health. Altogether,22 studies met the inclusion criteria, with a total sample size of 3409 employees post intervention, and 17 of the studies were included in the meta-analysis representing 20 interventions-control comparisons. The pooled results indicated small but positive overall effects of the interventions with respect to symptoms of depression and anxiety, but no effect on composite mental health measures. The interventions that included a direct focus on mental health had a comparable effect on depression and anxiety symptoms as did the interventions with an indirect focus on risk factors.\\

\noindent Jain et al. (2013) examined the burden of depression on work productivity in the US workforce. The Patient Health Questionnaire for depression severity, the Health and Work Performance Questionnaire and Work Productivity and Activity Impairment (WPAI) Questionnaire for absenteeism and presenteeism were used to survey 1051 full-time employees in February 2010. All levels of depression were associated with decreased work productivity as severity increased, presenteeism and absenteeism worsened.\\

\noindent In another study Serpa et al. (2014), investigated the major problems among Veterans specifically anxiety, depression and suicidal ideation. Pre- and Post-Mindfulness based Stress Reduction (MBSR) questionnaires were used to collect data on MBSR.s effectiveness among 79 veterans at an urban Veterans Health Administration medical facility and were used for comparisons among individuals. A meditation analysis was conducted to determine whether changes in mindfulness were related to changes in the other outcomes for a period of 9 weekly sessions. The results indicated significant reductions in anxiety, depressions, and suicidal ideation after MBSR training, improved mental health functioning scores although pain intensity and physical health functionality did not register improvements. Lister et al. (2023) conducted a qualitative study that investigated the  barriers and enablers to mental wellbeing and academic success that students experienced in distance learning. Narrative enquiry was used whereby students shared their own stories and tutors shared stories about the students they had supported.
In total, 60 themes were identified, consisting of 155 sub-themes and 783 coded
references. The barriers and enablers were then clustered into three overall categories;
study-related barriers/enablers, skills-related barriers/enablers and environmental barriers/enablers, and 10 sub-categories, curriculum, tuition, assessment, social skills, self-management, study skills, spaces, systems, people and life. It was found that barriers and enablers were existent in the majority of themes, implying that different aspects of study could be both barriers or enablers for mental wellbeing and study success, depending on who was experiencing them and how, for example, assessment feedback was a barrier for students when it was negative, but was an enabler when it was constructive and took place in an environment
of trust.\\

\noindent Another study by Magtibay et al. (2017) assessed the efficacy of blended learning to decrease stress and burnout among nurses through the use of the Stress Management and Resiliency Training (SMART) program. Consistent with blended learning, participants chose the format that met their learning styles and goals; Web-based, independent reading, facilitated discussions. The end points of mindfulness, resilience, anxiety, stress, happiness, and burnout were measured at baseline, postintervention, and 3-month follow-up to examine within-group differences. The results indicated statistically significant, clinically meaningful decreases in anxiety, stress, and burnout and increases in resilience, happiness, and mindfulness. \\

\noindent Levis et al. (2019), evaluates the accuracy of the Patient Health Questionnaire-9 (PHQ-9) for screening major depression in adults using individual participant data meta-analysis. The study, which was conducted by the DEPRESsion Screening Data (DEPRESSD) Collaboration, analyzed data from 58 studies, with a total of 17 357 participants. It was found that combined sensitivity and specificity was maximized at a cut-off score of 10 or above among studies using a semi structured interview. Secondly across cut-off scores of 5-15 ,sensitivity with semi structured interviews was 5\%-22\% higher than for fully structured interviews.\\

\noindent Leung et al. (2020) investigated the factor structure, reliability, and construct validity of the Patient Health Questionnaire-9 (PHQ-9) in a community sample of 10 933 Chinese adolescents. Confirmatory factor analysis (CFA) was used to examine the one-factor structure of the PHQ-9, and multiple-group CFAs were performed to test for measurement invariance across gender and age subgroups. Internal consistency reliability was evaluated using Cronbach alpha, and construct validity was assessed using Pearson correlation coefficients between PHQ-9 scores and anxiety, self esteem and perceived control. The results indicated that a one-factor model with three pairs of items correlations fitted the PHQ-9 data well and that measurement invariances by age and gender were supported. The PHQ-9 was also strongly correlated with anxiety, self esteem and perceived control.\\

\noindent In a more recent study, Levis et al. (2020), a meta-analysis of 44 studies was conducted to compare the sensitivity and specificity of the PHQ-2 alone and combined with PHQ-9. Individual participant data were obtained from 100 eligible studies comprising of 44 318 participants, of which 4572 had major depression.The combination of PHQ-2 followed by PHQ-9 had similar sensitivity but higher specificity compared with PHQ-9.\\

\noindent In summary, previous research works were conducted on mental health using Meta analysis, Qualitative methods, CFA, narrative enquiry (mathematical methods used)……….In this study factor analysis is used to…….(state what you want to do with FA)\\

\section*{REFERENCES}
\begin{enumerate}
\item Magtibay D.L., Chesak S.S., Coughlin K. and Sood A., (2017), Decreasing Stress and Burnout in Nurses, The Journal of Nursing Administration, 47(8), pp. 391-395.

\item Jain G., Roy A., Harikrishnan V., Yu S., Dabbous O. and Lawrence C., (2013), Patient-reported Depression Severity Measured by the PHQ-9 and Impact on Work Productivity, Journal of Occupational and Environmental Medicine, 55(3), pp. 252-258.

\item Serpa J.G., Taylor S.L. and Tillisch K., (2014), Mindfulness-based Stress Reduction (MBSR) Reduces Anxiety, Depression and Suicidal Ideation in Veterans, Medical Care, 52(12), pp. 19-24

\item Johnson R.A. and Wichern D.W., Applied Multivariate Statistical Analysis, Pearson Education, Inc., 2007

\item Leung D.Y.P., Chiang V.C.L., Mak Y.W., Loke A.Y. and Leung S.F., (2020), Measurement Invariances of the PHQ-9 across Gender and Age Groups in Chinese Adolescents, AsiaPacific Psychiatry,12, pp. 1-6.

\item Levis B., Benedetti A. and Thombs B.D., (2019), Accuracy of Patient Health Questionnaire9 (PHQ-9) for Screening to Detect Major Depression, British Medical Journal, 365, Pp. 1-10. 

\item Levis B., Sun Y., He C., Wu Y., Krishnan A., Bhandari P.M., Neupane D., Imran M., Brehaut E., Negeri Z., Fischer F.H., Benedetti A. and Thombs B.D., (2020), Accuracy of the PHQ-2 Alone and in Combination with the PHQ- 9 for Screening to Detect Major, Depression, Systematic Review and Meta Analysis, 323(22), pp .2290-2300 

\item Martin A., Sanderson K. and Cocker F., (2009), Meta Analysis of the Effects of Health Promotion Intervention in the Workplace on Depression and Anxiety symptoms, Scandinavian Journal of Work, 35 (1), pp. 7 – 18

\item Laing S.S. and Jones M.W., Anxiety and Depression Mediate the Relationship between Perceived Workplace Health Support and Presenteeism, 58(11), pp. 1144-1149.

\end{enumerate}
\end{document}
